\section{Details of TeGA}
\subsection{Details of the generated data by TeGA}
We provide samples of data generated by TeGA (\cref{fig:samples_passed} and \cref{fig:samples_filtered}).
The generated data which pass consistency filtering are shown in \cref{fig:samples_passed} and the generated data which are filtered out are shown in \cref{fig:samples_filtered}.
\cref{tab:class_ratios} compares some of the prompts used in the generation with the percentage that passed through the filter.

Classes that pass consistency filtering in large numbers tend to have clear class names. 
On the other hand, those that frequently fail consistency filtering often have ambiguous class names or simple shapes that make correspondence difficult to capture. 
For example, the 'birdhouse' has no fixed shape and could be a mesh or a wooden box. 
In such cases, it is difficult to generate and evaluate.
Conversely, the 'bicycle' has fixed the shape which is a vehicle with two wheels, a frame, handlebars for steering, and pedals for propulsion.
Some classes, despite having clear class names and not being simple in shape, are still prone to failing consistency filtering. 
% This is likely because these classes are not included in Point-E's training data.
\begin{table}[h!]
    \centering
    \adjustbox{width=\linewidth}{
    \begin{tabular}{@{}lcc@{}}
        \toprule
        \textbf{Class Name} & \textbf{Count (Filtered/Generated)} & \textbf{Percentage (\%)} \\ 
        \midrule
        airplane    & 3777/4045   & 93.38\%         \\ 
        bicycle     & 449/831     & 54.03\%         \\ 
        birdhouse   & 0/73        & 0.00\%          \\ 
        bookshelf   & 59/452      & 13.05\%         \\ 
        bottle      & 433/498     & 86.95\%         \\ 
        chair       & 6473/6778       & 95.05\%         \\ 
        clock       & 44/651      & 6.76\%         \\ 
        display     & 8/1093      & 7.32\%         \\ 
        lamp        & 1052/2318   & 45.38\%         \\ 
        telephone   & 42/1089     & 3.86\%         \\ 
        \bottomrule
    \end{tabular}}
    \caption{ShapeNet class name and the percentage that passed the filtering of the generation. A high percentage of the class names that are ambiguous are caught in the filtering process.}
    \label{tab:class_ratios}
\end{table}

\subsection{Details of the consistency filtering by TeGA}
\label{app:details_tega}
% \noindent{\textbf{consistency filtering by GPT-4.}}
In \cref{tab:prompt_examples_success} and \cref{tab:prompt_examples_failed_total_ok}, \cref{tab:prompt_examples_failed_total_ng}, we show an example of a part of TeGA using GPT-4.
This evaluation is instructed by pre-defined prompt.
The prompt is based on T3Bench~\cite{he2023t}.
In T3Bench, the prompt for generation assumes a concrete long sentence so there needs only to perform sense-level comparisons between caption and prompt.
However, in our evaluation, since the prompt for generation is the class name of ShapeNet, which is a word, T3Bench always returns low score.
For these reasons, it is necessary to account for conceptual similarities and we redesigned the prompt for evaluation.

In the generated captions, key information is sometimes placed at the end of the sentence, and low scores are output by GPT-4 in such cases.
% these sentense results in lower evaluation scores by GPT-4. 
For example, in the case of \cref{tab:prompt_examples_failed_total_ok}, despite the presence of the same classmate word 'sofa', a low score is output. In such cases, performing word-level matching outputs a high score and is able to complement GPT-4's shortcomings.

% \noindent{\textbf{consistency filtering by the comparison.}}
\addtocounter{table}{3}

\newcommand{\cmark}{\checkmark} % ✓
\newcommand{\xmark}{\textbf{\ding{55}}} % ✗

\begin{table}[t]
    \centering
    \adjustbox{width=0.8\linewidth}{
    \begin{tabular}{@{}ccc|c@{}}
        \toprule
        \multirow{2}*{Text-PC} & \multirow{2}*{Image-PC} & \multirow{2}*{Text-3D} & ModelNet-40 \\ 
        \cmidrule(l){4-4}
        & & & Top1 \\ 
        \midrule
        \cmark & \xmark & \xmark & 56.3 \\
        % & 52.9 & 78.7 & 86.5 \\ % 197
        % \xmark & \cmark & \cmark & \xmark & 38.4 & 38.5 & 61.6 & 73.4 \\ % 1 163
        \xmark & \cmark & \xmark & 35.9 \\
        % & 36.4 & 54.6 & 65.2 \\ % 2 131
        \xmark & \xmark & \cmark & 0.4 \\
        % & 5.0 & 7.6 & 12.5 \\ % 
        \xmark & \cmark & \cmark & 60.3 \\
        % & 60.6 & 79.3 & 87.1 \\ % 3 165
        \cmark & \cmark & \xmark & 42.2 \\
        % & 39.8 & 62.0 & 73.3 \\ % 4 144
        \cmark & \cmark & \cmark & \textbf{73.3} \\
        % & \textbf{68.7} & \textbf{88.8} & \textbf{93.6} \\ 
        \bottomrule
    \end{tabular}}
    \caption{Comparison of the final accuracy achieved when MixCon3D's contrastive learning is limited.  A check mark indicates an active contrastive learning setting, while a cross mark denotes an inactive setting. }
    \label{tab:multi-modal}
\end{table}

\addtocounter{table}{-4}
\addtocounter{figure}{2}
\begin{figure}[ht!]
    \centering
    \includegraphics[width=\linewidth]{figure/shapenet_vs_synth2.pdf}
    \caption{Visualization of features. Features with the same class name are clustered close to each other. However, real data and synthetic data are not clustered close to each other.}
    \label{fig:visualize_figure}
\end{figure}
\addtocounter{figure}{-3}

\section{Additional Experimental Results}
\noindent{\textbf{Multi-modal Pretraining.}}
We investigate which modality interactions contribute to training for language-image-3D tasks using synthetic data. 
Specifically, it compares the final accuracy achieved when MixCon3D's contrastive learning is applied to only a subset of the modalities. \cref{tab:multi-modal} presents the results. 
A check mark indicates an active contrastive learning setting, while a cross mark denotes an inactive setting. 
Note that the text-image contrastive learning is always active due to knowledge distillation from OpenCLIP~\cite{cherti2023reproducible}.

Surprisingly, high performance is achieved even without utilizing text-point cloud contrastive learning. Additionally, although training succeeds without text-3D contrastive learning, the accuracy is higher when it is included, indicating that text-3D contrastive learning definitely contributes to improving performance.

\begin{figure*}
    \centering
    \includegraphics[width=1.0\linewidth]{figure/samples_passed2_reduced.pdf}
    \caption{Samples of data generated by TeGA.}
    \label{fig:samples_passed}
\end{figure*}

\begin{figure*}
    \centering
    \includegraphics[width=1.0\linewidth]{figure/samples_filtered.pdf}
    \caption{Samples of data generated but filtered out in TeGA.}
    \label{fig:samples_filtered}
\end{figure*}

\noindent{\textbf{Visualize Feature.}}
\cref{fig:visualize_figure} shows the visualization of the predicted feature by MixCon3D.
The features of the three classes—airplane, table, and sofa—are extracted from both ShapeNet data and synthetic data, and dimensionality reduction is performed using t-SNE.
Surprisingly, not only are different classes separated, but even the ShapeNet data and synthetic data of the same class are distinctly separated.

\begin{table*}[ht] % table* は twocolumn 用
    \centering
    \begin{tabular}{@{}p{2.5cm}p{12cm}@{}}
        \toprule
        \textbf{Prompt} & 
         You are an assessment expert responsible for prompt-prediction pairs. Your task is to score the prediction according to the following requirements:\newline
         1. Evaluate the recall, or how well the prediction covers the information in the prompt. If the prediction contains information that does not appear in the prompt, it should not be considered as bad.\newline
        2. Assign a score between 1 and 5, with 5 being the highest. Do not provide a complete answer; give the score in the format: 3\newline
        3. add points if the prediction and prompt are conceptually close (e.g. similar in appearance). (e.g., bike and bycicle and table and chair are close)\newline
        4. since the prompt is at the word level, it is inevitable that some detailed information is missing, so exclude it from the point deduction. \newline \newline
        prompt: car\newline
        prediction: A 3D rendering of a car with a pink and white exterior and a pink interior with red streaks\\
        \midrule
        \textbf{Answer:} &
        5 \\
        \midrule
        \textbf{Word matching} &
        5 \\
        \midrule
        \textbf{Total} &
        10 (safe) \\
        % \midrule
        % & Reply with the format of ``index\#class\#short reason (no more than 10 words)". \\
        % \midrule
        % \textbf{Examples:} & \\
        % Input: This is a 3D object model of a cartoon white truck. \newline
        % Output: 7\#car\#Closest match to "car" in categories. \newline

        % Input: A green leaf in a flower pot. \newline
        % Output: 26\#plant\#The primary subject "leaf" directly indicates a plant. \newline

        % Input: It’s difficult to determine the exact type of this object due to insufficient details. But it seems to be like a piece of furniture. \newline
        % Output: 33\#table\#Randomly select one kind of furniture from the list. \newline

        % Input: I cannot determine the specific type of the object without additional information or context. \newline
        % Output: -1\#NA\#Cannot infer. \\
        % \midrule
        % & Now analyze the following: \newline
        % Input: \{model\_output\} \newline
        % Output: \\
        \bottomrule
    \end{tabular}
    \caption{The process of consistency filtering in TeGA. Captions and prompts are evaluated by GPT-4 and by word-matching. Both of evaluation are judged to be high.}
    \label{tab:prompt_examples_success}
\end{table*}
\begin{table*}[ht] % table* は twocolumn 用
    \centering
    \begin{tabular}{@{}p{2.5cm}p{12cm}@{}}
        \toprule
        \textbf{Prompt} & 
         You are an assessment expert responsible for prompt-prediction pairs. Your task is to score the prediction according to the following requirements:\newline
         1. Evaluate the recall, or how well the prediction covers the information in the prompt. If the prediction contains information that does not appear in the prompt, it should not be considered as bad.\newline
        2. Assign a score between 1 and 5, with 5 being the highest. Do not provide a complete answer; give the score in the format: 3\newline
        3. add points if the prediction and prompt are conceptually close (e.g. similar in appearance). (e.g., bike and bycicle and table and chair are close)\newline
        4. since the prompt is at the word level, it is inevitable that some detailed information is missing, so exclude it from the point deduction. \newline \newline
        prompt: sofa\newline
        prediction: A modern, cream-colored sofa\\
        \midrule
        \textbf{GPT-4 Answer} &
        1 \\
        \midrule
        \textbf{Word matching} &
        5 \\
        \midrule
        \textbf{Total} &
        6 (safe) \\
        % \midrule
        % & Reply with the format of ``index\#class\#short reason (no more than 10 words)". \\
        % \midrule
        % \textbf{Examples:} & \\
        % Input: This is a 3D object model of a cartoon white truck. \newline
        % Output: 7\#car\#Closest match to "car" in categories. \newline

        % Input: A green leaf in a flower pot. \newline
        % Output: 26\#plant\#The primary subject "leaf" directly indicates a plant. \newline

        % Input: It’s difficult to determine the exact type of this object due to insufficient details. But it seems to be like a piece of furniture. \newline
        % Output: 33\#table\#Randomly select one kind of furniture from the list. \newline

        % Input: I cannot determine the specific type of the object without additional information or context. \newline
        % Output: -1\#NA\#Cannot infer. \\
        % \midrule
        % & Now analyze the following: \newline
        % Input: \{model\_output\} \newline
        % Output: \\
        \bottomrule
    \end{tabular}
    \caption{The process of consistency filtering in TeGA. Captions and prompts are evaluated by GPT-4 and by word-matching. The prompts appear in the caption, but are judged low by GPT-4, while word-matching is judged high. On the other hand, the word-matching score is high.}
    \label{tab:prompt_examples_failed_total_ok}
\end{table*}


% \begin{table*}[ht]
%     \centering
%     % 1つ目の表
%     \begin{minipage}{0.3\linewidth}
%         \centering
%         \begin{tabular}{@{}p{1cm}p{6cm}@{}}
%             \toprule
%             \textbf{Prompt} & prompt: sofa \newline
%             prediction: A modern, cream-colored sofa\\
%             \midrule
%             \textbf{GPT-4 Answer} &
%             1 \\
%             \midrule
%             \textbf{Word matching} &
%             5 \\
%             \midrule
%             \textbf{Total} &
%             6 (safe) \\
%             \bottomrule
%         \end{tabular}
%         \caption{Table 1}
%         \label{tab:table1}
%     \end{minipage}
%     \hfill
%     % 2つ目の表
%     \begin{minipage}{0.3\linewidth}
%         \centering
%         \begin{tabular}{@{}cc@{}}
%             \toprule
%             X & Y \\
%             \midrule
%             5 & 6 \\
%             7 & 8 \\
%             \bottomrule
%         \end{tabular}
%         \caption{Table 2}
%         \label{tab:table2}
%     \end{minipage}
%     \hfill
%     % 3つ目の表
%     \begin{minipage}{0.3\linewidth}
%         \centering
%         \begin{tabular}{@{}cc@{}}
%             \toprule
%             P & Q \\
%             \midrule
%             9  & 10 \\
%             11 & 12 \\
%             \bottomrule
%         \end{tabular}
%         \caption{Table 3}
%         \label{tab:table3}
%     \end{minipage}
% \end{table*}
\input{tables/tab_prompt_example_failed_total_ng}